% =============================================================================
\subsection{Synthesis of Key Findings}
% =============================================================================

The three studies presented in Section~\ref{sec:results} collectively establish a nuanced picture: quantum \textit{improvement} is not universal but is \textit{formulation-dependent} and \textit{scale-dependent}. We discuss each study's contribution in turn, followed by an analysis of hardware considerations and future directions.

% -----------------------------------------------------------------------------
% FIRST HALF: Study-by-study conclusions
% -----------------------------------------------------------------------------

\subsubsection{Study~1: Hybrid Solvers --- Performance Without Transparency}

Study~1.A (Section~\ref{subsec:hybrid_benchmarks}) established that classical Gurobi consistently finds optimal solutions in under 1.2~seconds across all problem scales (10 to 1,000 patches), confirming its dominance on structured MILP formulations. D-Wave's LeapHybridCQMSolver, by comparison, maintains a constant time profile of approximately 5.3 to 5.5~seconds but trails Gurobi in solution quality: the hybrid CQM's optimality gap ranges from 10.5\% at 10~farms down to 0.8\% at 50~farms (Table~\ref{tab:hybrid_cqm_performance}). At larger scales (200+ patches), the hybrid CQM solver becomes infeasible, producing solutions with coverage exceeding available land by 15--17$\times$ (Table~\ref{tab:hybrid_cqm_patch}). This indicates that the hybrid solver's internal heuristics prioritize objective optimization over strict constraint satisfaction when problem density increases --- a trade-off that is invisible to the user and contrasts with Gurobi's consistently feasible solutions.

Beyond the quality gap, the most consequential finding from Study~1.A is the minimal QPU contribution within the hybrid solver. Post-hoc analysis revealed that actual QPU annealing time is consistently 35--70\,ms, constituting only \textbf{1.3\%} of total wall-clock time (Figure~\ref{fig:hybrid_solver_performance}). The remaining 98.7\% is classical preprocessing --- problem decomposition, embedding search, and solution refinement. This opacity motivated the development of transparent decomposition methods in Study~2, where quantum and classical contributions are explicitly separated and measured.

Study~1.B demonstrated that the LeapHybridBQMSolver (Table~\ref{tab:dwave_bqm_scaling}) successfully solved QUBO-encoded problems across all scales (up to 55,310 variables), whereas classical Gurobi consistently timed out or returned infeasible solutions on the same formulation (Table~\ref{tab:gurobi_qubo_degradation}). This early indication that problem encoding determines solver effectiveness foreshadows the formulation-dependent quantum \textit{improvement} observed in Study~3.

\subsubsection{Study~2: Pure QPU Decomposition --- Transparent Scaling}

Study~2 (Section~\ref{subsec:qpu_decomposition}) addressed the black-box limitation by developing eight transparent decomposition strategies. The central finding is that \textbf{pure QPU annealing time scales linearly} with problem size: $T_{\text{QPU}} = O(f)$, where $f$ is the number of farms. As shown in Table~\ref{tab:qpu_time_scaling}, pure QPU time remains under 30~seconds even at 1,000~farms (27,027 variables), while classical embedding consumes 95--99\% of total runtime at scale --- growing from 1.2\,s at 10~farms to 3,473.6\,s at 1,000~farms.

The solution quality comparison at 1,000~farms (Table~\ref{tab:quality_1000farms}) reveals a quality-feasibility-diversity trade-off across decomposition strategies. Coordinated decomposition achieves the best objective (31.8\% gap) but introduces 23~constraint violations. Multilevel(10) achieves 39.9\% gap with zero violations and uses all 27~crops --- maximum diversity. By contrast, Gurobi allocates 99.6\% of land to spinach, satisfying diversity constraints minimally.

The violation healing analysis (Section~\ref{subsec:qpu_decomposition}) quantified that constraint violations in Study~2 explain only approximately 7\% of the optimality gap, with the remaining 93\% attributable to decomposition boundary effects and stochastic sampling variance. Methods such as CQM-First PlotBased and Louvain maintain 100\% feasibility across all tested scales (Table~\ref{tab:violation_by_strategy}), demonstrating that zero-violation quantum solutions are achievable. The comprehensive benchmark results (Figures~\ref{fig:qpu_benchmark_small},~\ref{fig:qpu_benchmark_large}, and~\ref{fig:qpu_benchmark_comprehensive}) establish that Multilevel(10) and HybridGrid(10,9) achieve the best balance of solution quality, computational efficiency, and constraint satisfaction across all scales.

\subsubsection{Study~3: Variant~B --- Quantum \textit{Improvement} on Crop Rotation}

Study~3 (Section~\ref{subsec:quantum_advantage}) provides the strongest evidence for practical quantum \textit{improvement}. Across 13~rotation scenarios, the QPU achieves \textbf{3.80$\times$ higher benefit values} than Gurobi on average (Table~\ref{tab:qpu_advantage}), with the \textit{improvement} ratio increasing from 2.51$\times$ at 180~variables to 5.35$\times$ at 16,200~variables. This scale-dependent growth is a central result: quantum \textit{improvement} does not merely persist at larger scales but \textit{intensifies}.

The underlying mechanism is Gurobi's inability to solve the crop rotation MIQP formulation effectively. As documented in Table~\ref{tab:gurobi_struggles}, Gurobi times out on 11 of 13~scenarios with an average MIP gap of 16,308\%. For the 6-Food (large) formulation, MIP gaps reach 352,822\%, indicating that Gurobi's best feasible solution after 300~seconds is astronomically far from the proven lower bound. The QPU, by contrast, explores solution regions that Gurobi's branch-and-bound algorithm cannot reach within practical time limits.

The timing analysis (Table~\ref{tab:timing_comparison}) confirms that the embedding bottleneck pattern observed in Study~2 persists in Study~3: pure QPU time totals only 10.88~seconds across all 13~scenarios (1.1\% of wall time), scaling linearly at $T = 0.78 \cdot N_{\text{vars}} + 51$\,ms. Classical embedding and coordination consume the remaining 99\% of runtime. The formulation-specific analysis (Figure~\ref{fig:split_analysis}) further reveals that this overhead scales with \textit{farm count} rather than variable count, as discussed in Section~\ref{sec:formulation_decomposition_structure}.

\subsubsection{Cross-Study Conclusions}

Three patterns emerge consistently across all studies:

\begin{enumerate}
    \item \textbf{Quantum \textit{improvement} is formulation-dependent.} Classical Gurobi dominates on Variant~A's structured MILP (Study~1.A), but the QPU achieves 3.80$\times$ higher benefit on Variant~B's MIQP rotation formulation (Study~3). Quantum \textit{improvement} emerges specifically for problems where quadratic coupling creates frustrated landscapes~\cite{barahona1982computational, lucasIsingFormulationsMany2014} that defeat classical branch-and-bound.

    \item \textbf{Classical embedding is the dominant bottleneck.} In Study~1.A, QPU time is 1.3\% of hybrid wall-clock time. In Study~2, embedding consumes 95--99\% of total time at 1,000~farms (Table~\ref{tab:qpu_time_scaling}). In Study~3, pure QPU time is 1.1\% across all 13~scenarios (Table~\ref{tab:timing_comparison}). The pure QPU scaling ($O(f)$, linear) is already favorable; the bottleneck is the classical embedding step, which grows as $\sim O(f^{1.5})$~\cite{zbinden2020embedding}.

    \item \textbf{Decomposition methods produce practically valuable solutions.} Quantum decomposition yields diverse crop allocations (all 27~crops in Multilevel(10)) compared to Gurobi's monoculture optimum (99.6\% spinach, Table~\ref{tab:quality_1000farms}). For agricultural planning, this diversity provides resilience, nutritional variety, and soil health benefits that raw objective value does not capture~\cite{lin2011resilience, bowles2020longterm}. Where constraint violations occur (Study~2: Table~\ref{tab:violation_by_strategy}; Study~3: Table~\ref{tab:violation_impact}), their nature --- fields left temporarily fallow --- is easily corrected in post-processing and may even be agronomically desirable~\cite{dury2012cropping, hobbs2008conservation}.
\end{enumerate}

% -----------------------------------------------------------------------------
% SECOND HALF: Hardware considerations, noise, future directions, FTQC
% -----------------------------------------------------------------------------

% =============================================================================
\subsection{Hardware Considerations and Future Directions}
\label{subsec:hardware_future}
% =============================================================================

\subsubsection{QPU Noise and Chain Break Management}

The D-Wave Advantage system operates as a noisy intermediate-scale quantum (NISQ) device~\cite{preskillQuantumComputingNISQ2018, boothby2020next}, subject to thermal fluctuations, control errors, and inter-qubit coupling imperfections. The primary hardware artifact in our experiments was \textbf{chain breaks}: when a logical variable requires multiple physical qubits (a chain), thermal noise can cause disagreement among chained qubits~\cite{choi2008minor}.

Our decomposition strategy was explicitly designed to minimize this effect. Farm-level subproblems (27~variables) achieve chain lengths $\leq 1.2$ on the Pegasus topology~\cite{boothby2020next}, and native clique embeddings (15--20 qubits fully connected) exhibited zero chain breaks. Across all experiments, the chain break rate remained below 2\%, with auto-scaled chain strength (1.2--1.8$\times$ the maximum quadratic coefficient) proving consistently effective~\cite{venturelli2015quantum}. Automatic postprocessing via greedy descent~\cite{cai2014practical} adds negligible time ($<$100\,ms per sample) and repairs any remaining chain break artifacts.

The observable effect of noise was sample variance: the 100~samples per QPU call exhibited a distribution of objective values rather than converging to a single solution. This is expected behavior for quantum annealing --- the sampler explores a neighborhood around the minimum-energy state --- and is addressed by selecting the best sample from each batch. For future work, adaptive annealing schedules~\cite{marshall2019power, king2022coherent} and reverse annealing~\cite{venturelli2019reverse} --- starting from a classical initial solution and refining via quantum fluctuations --- could improve sampling of low-energy states in the frustrated rotation landscape.

\subsubsection{The Embedding Bottleneck}

The single most important finding across all three studies is that \textbf{classical embedding dominates end-to-end runtime}. Table~\ref{tab:qpu_time_scaling} quantifies this precisely: at 1,000~farms, embedding consumes 3,473.6\,s (99.4\% of total time) while pure QPU annealing requires only 21.78\,s (0.6\%). This pattern is consistent across Studies~2 and~3, and across both the Variant~A decomposition benchmarks and the Variant~B rotation scenarios.

Embedding complexity depends on problem connectivity and QPU topology~\cite{choi2008minor, zbinden2020embedding}. Our decomposition strategies explicitly control this trade-off:
\begin{itemize}
    \item \textbf{PlotBased:} Sparse per-farm subproblems (27 variables) $\to$ embedding in $<$0.5\,s per partition
    \item \textbf{Multilevel(10):} Medium subproblems ($\sim$270 variables) $\to$ embedding in 5--30\,s per partition
    \item \textbf{Direct QPU:} Full problem (27,027 variables) $\to$ embedding failure (Table~\ref{tab:quality_1000farms})
\end{itemize}

This hierarchy establishes a design principle: decomposition strategies should create subproblems that match hardware capabilities. For the current Pegasus topology, targeting 20--50 variable partitions with sparse connectivity ensures fast, reliable embedding. Any hardware improvement that reduces embedding overhead --- whether through higher connectivity, larger native cliques, or dedicated embedding co-processors~\cite{boothby2021zephyr} --- would have an outsized impact on end-to-end performance.

If embedding overhead were eliminated entirely, the 1,000-farm problem would complete in $\sim$22~seconds of pure QPU time instead of $\sim$3,495~seconds. This 159$\times$ acceleration --- achievable through hardware improvements alone, without algorithmic changes --- would place QPU performance within the same order of magnitude as Gurobi's 0.32-second solve time on Variant~A, and would dramatically outperform Gurobi's 300-second timeouts on Variant~B. For the Variant~B rotation scenarios, the QPU would deliver 3.80$\times$ higher benefit (Table~\ref{tab:qpu_advantage}) in under 11~seconds of total compute time --- a regime where no classical solver currently competes.

\subsubsection{Impact of Higher Qubit Connectivity}

The most directly relevant hardware improvement for our application is \textbf{increased qubit connectivity}. The current D-Wave Advantage system uses the Pegasus topology with degree-15 connectivity, supporting native cliques of 15--20 qubits~\cite{boothby2020next}. Our farm-level subproblems (27~variables for 27-Food, 18~variables for 6-Food) can embed within or near these native cliques, but larger subproblems (e.g., the 270-variable partitions in Multilevel(10)) require chains and multi-second embedding searches~\cite{cai2014practical}.

D-Wave's next-generation \textbf{Zephyr topology} offers degree-20+ connectivity, supporting native cliques of approximately 30--40 qubits~\cite{boothby2021zephyr}. For our problem, this advancement would have three concrete effects:

\begin{enumerate}
    \item \textbf{Larger chain-free subproblems:} Farm clusters of 5--7~farms (up to $\sim$126 variables for 6-Food, $\sim$189 for 27-Food) could embed without chains, reducing the number of partitions by 5--7$\times$ and proportionally reducing both QPU calls and coordination overhead.

    \item \textbf{Faster embedding for remaining chains:} Higher connectivity reduces the search space for MinorMiner~\cite{cai2014practical}, decreasing per-partition embedding time from the current 5--30\,s to projections of $<$1\,s for moderate subproblems.

    \item \textbf{Reduced boundary approximation error:} Fewer, larger partitions mean fewer inter-partition boundaries, reducing the decomposition approximation that accounts for $\sim$93\% of the current optimality gap. This could close the 30--40\% optimality gaps observed in Study~2 to substantially lower values.
\end{enumerate}

Combining these effects, we conservatively estimate that Zephyr-class connectivity could yield a 5--10$\times$ end-to-end speedup for the decomposition methods benchmarked in Study~2, and could reduce optimality gaps by 30--50\% through improved boundary coordination.

\subsubsection{Impact of Larger QPU Processors and Lower Noise}

Beyond connectivity, future QPUs with more qubits would enable embedding larger subproblems --- or even entire problems --- without decomposition. Current Pegasus systems have $\sim$5,760 qubits~\cite{boothby2020next}; next-generation systems may offer 10,000--20,000+ qubits~\cite{boothby2021zephyr}. For our 200-farm 27-Food rotation problem (16,200 variables), a QPU with sufficient qubits and connectivity could eliminate decomposition entirely, removing boundary approximation errors, solving the global optimization in a single annealing call ($\sim$20\,ms), and producing solutions that respect global constraints natively.

Improved qubit coherence and reduced thermal noise would further benefit our application. Our current chain break rate of $<$2\% is already low, but further reduction to $<$0.1\% would enable more aggressive penalty tuning. Currently, chain strength must be set conservatively (1.2--1.8$\times$ max coefficient) to prevent breaks, which can distort the energy landscape. Lower noise would allow tighter chain strengths, producing more faithful representations of the QUBO objective and potentially improving solution quality. Additionally, reduced sample variance would increase the probability that the best sample from each batch is close to the true ground state --- particularly impactful for problems with many nearly-degenerate local minima, such as our frustrated rotation landscapes.

\subsubsection{Parallelization and Multi-QPU Architectures}

Our decomposition approach is inherently parallelizable: independent farm clusters can be solved simultaneously on separate QPUs. For the 1,000-farm Multilevel(10) benchmark (102~partitions), a hypothetical 10-QPU array would reduce pure QPU time from 21.78\,s to $\sim$2.2\,s. The embedding bottleneck would also parallelize: each QPU independently embeds its assigned partitions, reducing embedding wall-clock time by a factor of 10.

Under such a parallel architecture, the 1,000-farm problem currently requiring $\sim$3,495\,s would complete in approximately $\sim$350\,s --- comparable to a single Gurobi timeout, but delivering a solution rather than hitting a time limit. Combined with Zephyr-class connectivity improvements, parallel multi-QPU execution could bring the 1,000-farm Variant~A problem to completion in under 100~seconds, and could solve the Variant~B rotation scenarios with 3.80$\times$ higher benefit in comparable time.

\subsubsection{Conditions for Realistic Quantum \textit{Improvement}}

Based on the evidence from all three studies, we identify the following conditions under which practical quantum \textit{improvement} could transition from the demonstration presented here to a deployment-ready capability:

\begin{enumerate}
    \item \textbf{Problem structure must favor QUBO encoding.} Quantum \textit{improvement} emerges specifically for problems where quadratic coupling (rotation constraints, synergy terms) creates frustrated landscapes~\cite{barahona1982computational, lucasIsingFormulationsMany2014} that defeat classical branch-and-bound. On Variant~A's structured MILP, Gurobi solves in $<$1.2\,s at all scales (Section~\ref{subsec:hybrid_benchmarks}); on Variant~B's MIQP rotation formulation, Gurobi times out on 11/13~scenarios (Table~\ref{tab:gurobi_struggles}).

    \item \textbf{Classical embedding overhead must be reduced.} Hardware improvements in connectivity (Zephyr and beyond)~\cite{boothby2021zephyr}, algorithmic improvements to MinorMiner~\cite{cai2014practical}, or the development of topology-aware problem formulations could each contribute to reducing the embedding overhead that currently consumes 95--99\% of runtime (Table~\ref{tab:qpu_time_scaling})~\cite{zbinden2020embedding}.

    \item \textbf{Decomposition quality must improve.} The 30--40\% optimality gaps observed in Study~2 are acceptable for many practical applications but leave room for improvement. Larger native cliques, better boundary coordination algorithms, and hardware-aware decomposition strategies co-designed with QPU topology could close this gap substantially.

    \item \textbf{Scale must be sufficient.} The QPU \textit{improvement} ratio grows with problem size (from 2.51$\times$ at 180~variables to 5.35$\times$ at 16,200~variables in Study~3). The linear QPU time extrapolation ($T = 0.78 \cdot N_{\text{vars}} + 51$\,ms) predicts $\sim$78~seconds for a 100,000-variable problem. At scales relevant to regional agricultural planning (1,000--10,000~farms), quantum \textit{improvement} is projected to be most pronounced.
\end{enumerate}

These conditions are not speculative: they are directly derived from the quantitative findings of the three studies. The gap between current demonstration and deployment-ready \textit{improvement} is primarily an engineering challenge (embedding overhead, QPU connectivity) rather than a fundamental theoretical barrier.

\subsubsection{Fault-Tolerant Quantum Computing Considerations}

While our work uses quantum annealing rather than gate-based quantum computing, the broader trajectory toward fault-tolerant quantum computation (FTQC)~\cite{campbell2017roads} has implications for the problem class studied here. Fault-tolerant quantum optimization algorithms (e.g., quantum approximate optimization algorithm variants~\cite{edwardfarhiQuantumApproximateOptimization2014} with error correction) could in principle solve QUBO problems with guaranteed approximation ratios, removing the stochastic variability inherent in current quantum annealing.

However, for the near term, the relevant question is whether \textit{current NISQ-era} quantum annealers can deliver practical \textit{improvement} for our specific problem class. Our results suggest the answer is conditionally affirmative: the QPU delivers 3.80$\times$ higher benefit on Variant~B (Table~\ref{tab:qpu_advantage}) with the \textit{improvement} ratio growing with scale, but end-to-end wall-clock time is currently dominated by classical embedding overhead rather than quantum computation limits.

% =============================================================================
\subsection{Conclusions}
% =============================================================================

Phase~3 established quantum annealing as a viable computational approach for large-scale agricultural optimization, with demonstrated quantum \textit{improvement} on the crop rotation formulation and clear pathways to enhanced performance through hardware evolution. The key conclusions, each grounded in specific experimental evidence, are:

\begin{enumerate}
    \item \textbf{Quantum \textit{improvement} is formulation-dependent.} Classical Gurobi solves Variant~A's structured MILP in under 1.2~seconds at all scales (Table~\ref{tab:hybrid_cqm_performance}), but times out on 11/13 Variant~B rotation scenarios with average MIP gaps of 16,308\% (Table~\ref{tab:gurobi_struggles}). The QPU achieves 3.80$\times$ higher benefit on Variant~B (Table~\ref{tab:qpu_advantage}), demonstrating that problem encoding determines which computational paradigm succeeds.

    \item \textbf{Quantum \textit{improvement} grows with scale.} The QPU benefit ratio increases from 2.51$\times$ at 180~variables to 5.35$\times$ at 16,200~variables (Table~\ref{tab:qpu_advantage}), and pure QPU time scales linearly at 0.78\,ms/variable (Table~\ref{tab:timing_comparison}). This scaling behavior, combined with the expected superlinear growth of classical solver difficulty, projects increasing quantum \textit{improvement} at problem sizes relevant to regional agricultural planning.

    \item \textbf{Classical embedding is the primary bottleneck.} Pure QPU time constitutes only 0.6--14.9\% of total runtime (Table~\ref{tab:qpu_time_scaling}), with embedding search consuming the remainder. This finding identifies the principal target for both hardware improvements (higher connectivity) and algorithmic optimization (topology-aware decomposition).

    \item \textbf{Transparency reveals genuine quantum contribution.} The transition from black-box hybrid solvers (Study~1.A) to transparent decomposition methods (Study~2) revealed that hybrid solver QPU usage is only 35--70\,ms (1.3\% of wall-clock time), motivating explicit quantum-classical separation. This transparency is essential for honest assessment of quantum contribution and for projecting performance on future hardware.

    \item \textbf{Quantum solutions offer practical diversity benefits.} Decomposition methods naturally produce diverse crop allocations (all 27~crops in Multilevel(10)) compared to Gurobi's monoculture optimum (99.6\% spinach), as documented in Table~\ref{tab:quality_1000farms}. For agricultural planning, this diversity provides resilience, nutritional variety, and soil health benefits that raw objective value does not capture~\cite{lin2011resilience, bowles2020longterm}.

    \item \textbf{Violations are a manageable trade-off.} The 24.2\% violation rate in Study~3 (Table~\ref{tab:violation_impact}) explains only $\sim$7\% of the objective gap (Figure~\ref{fig:violation_impact}); zero-violation methods exist for Study~2 (Table~\ref{tab:violation_by_strategy}); and the nature of violations (fields left fallow) is easily corrected in post-processing.

    \item \textbf{Future hardware improvements target the right bottleneck.} Higher qubit connectivity (Zephyr topology), larger processors, lower noise, and multi-QPU parallelization all directly address the classical embedding overhead that currently dominates runtime. If embedding overhead were eliminated, the 1,000-farm problem would complete in $\sim$22~seconds of pure QPU time --- a regime where quantum \textit{improvement} is decisive.
\end{enumerate}

These results justify progression to Phase~4 QPU-based proof-of-concept implementation, focusing on real-world deployment scenarios with multi-period rotation planning for agricultural systems. This Phase~3 study demonstrates practical quantum \textit{improvement} for a real-world optimization problem, provides transparent accounting of quantum versus classical computation, and establishes that quantum \textit{improvement} grows at larger scales. These findings provide a strong foundation for Phase~4 deployment, where the quantum optimization approach will be tested with real farm cooperatives in the field.


